\section{The arithmetic projection theorem}\label{sec:9}
\begin{definition}[\texttt{SemanticQuadraticGeneral/Defs.\allowbreak{}lean}]\label[definition]{definition:9.1}
For a form $H$ in $r+1$ variables write $\mathrm{fiberSetG}\ H\ ℓ\ q := \{c ∈ ℕ \mid ∃ u ∈ ℤ^{r}, H(c,u) + ℓ·(c,u) + q = 0\}$ and $\mathrm{fboundG}\ K\ ℓ\ q := K·((1 + Σ|ℓ_{i}|)^{2} + |q|)$. Then $\mathrm{DichotomyG}\ H :⟺ ∃ K$, for every affine datum $(ℓ, q)$, $\mathrm{fiberSetG}\ H\ ℓ\ q$ is infinite or bounded by $\mathrm{fboundG}\ K\ ℓ\ q$. $\mathrm{QuadraticProjectionDichotomyZ}\ r :⟺$ every homogeneous quadratic $H$ in $r+1$ variables satisfies $\mathrm{DichotomyG}\ H$.

\end{definition}
\begin{theorem}[\texttt{quadraticProjectionDichotomyZ\_\allowbreak{}all}, unconditional]\label[theorem]{theorem:9.2}
For every $r$, $\mathrm{QuadraticProjectionDichotomyZ}\ r$.

\end{theorem}
\begin{proof}
This is what \cref{sec:F} proves, in full and without appeal to the formalization: Part I settles the full-rank case for $r ≥ 2$, Part II the singular case with its dispatch and the size transport back to the original affine data, Part V the small-arity and two-variable leaves the dispatch calls, and Part IV assembles them.

\end{proof}
\begin{theorem}[data-carrying form; \texttt{SemanticQuadraticGeneral/TopDichotomyData.\allowbreak{}lean}, \texttt{quadraticProjectionDichotomyZ\_\allowbreak{}all\_\allowbreak{}data}, \texttt{quadraticProjectionDichotomyZ\_\allowbreak{}structured}, unconditional]\label[theorem]{theorem:9.3}
The same statement with "infinite" replaced by "carries a structured family": $∃ K$, for every $(ℓ, q)$, either there is a \texttt{StructuredFamily (fiberSetG H ℓ q)}, or every element is bounded by $\mathrm{fboundG}\ K\ ℓ\ q$. \cref{theorem:9.2} follows by forgetting the family (\texttt{dichotomyG\_\allowbreak{}of\_\allowbreak{}dichotomyGData}).

\end{theorem}
\begin{proof}
The complete proof of \cref{theorem:9.2,theorem:9.3} is \cref{sec:F} of this paper, where it is given as twenty-five numbered lemmas with their proofs; the formal development verifies the same argument, and \cref{sec:12} records the correspondence. What follows in this section is a road map of that proof, so that a reader can see the shape of the argument before reading it in full. \cref{theorem:2.1,theorem:8.5} both rest on this one theorem, and a reader is entitled to know what kind of argument carries that weight.

\emph{Normal form.} The quadratic form H is presented by its Gram data, and one first brings it into a shape in which the radical is visible. The reduction uses a unimodular \emph{congruence} $U^{T}$ G U — not a Smith-normal-form left–right equivalence, which would destroy the quadratic form — obtained from the saturation of the integer kernel of G together with a complement of the torsion-free quotient (\texttt{RadicalDecomposition.\allowbreak{}lean}, \texttt{Unimodular3.\allowbreak{}lean}). This splits the problem into a full-rank part and a singular part, and the affine data transform correctly under the congruence, which is what makes the bound fboundG survive the change of coordinates.

\emph{Full-rank case.} The fibre equation is a nondegenerate quadratic condition, and the dichotomy is proved at each fibre point rather than for the fibre as a whole. The top-level split is on the sign of α(2αc₀+β): non-positive forces $c_{0} ≤ |β|$ and the point is bounded; positive, together with a signed non-square direction, produces an infinite orbit (Part I §§3–5). The case $α = 0$ is separate (§6), and the direction is produced by a sign trichotomy of the adjugate form, with the semidefinite alternative $α·\tilde{W} ≥ 0$ everywhere treated on its own (§§7–8). The classical isotropic–anisotropic distinction enters inside that last step rather than at the top: an isotropic form yields a line of solutions and a polynomial progression (Lemma V.3); an anisotropic one brings in the unit group of a real quadratic order (Lemma V.2), and a suitable initial solution of the shifted norm equation generates an infinite orbit under a unit preserving the congruence conditions the branch requires. (Nothing below needs the stronger classical statement that all solutions fall into finitely many such orbits; one infinite orbit suffices.) Iterating the unit produces the Pell-unit orbit of \cref{definition:8.2}(2), with the recurrence read off from multiplication by x₀ + y₀√D; the affine projection with denominator appears because the output coordinate is recovered from the orbit point only after clearing a fixed denominator (\texttt{FullRankAssemblyData.\allowbreak{}lean}).

\emph{Singular case.} The form degenerates along its radical, and the reduction does \textbf{not} proceed by an induction on the number of unknowns. Part II dispatches on the rank k₀ of the nondegenerate quotient: $k_{0} = 0$ to Lemma $4, k_{0} = 1$ to Lemma $6, k_{0} ≥ 2$ to Lemma 5 (\texttt{QuotientRankZero.\allowbreak{}lean}, \texttt{QuotientRankOne.\allowbreak{}lean}). Underneath the dispatch sit the radical affine coefficients (\texttt{RadicalAffine.\allowbreak{}lean}), a nonzero-minor periodicity argument, a proportional compression, and the size transport back to the original affine data. The rank-zero leaf is the largest single case: all three of its infinite branches turn out to be progressions (an identity family, a periodic family, and a square progression).

\emph{The moving-modulus branch.} One branch of the singular case resists both treatments: the fibre condition becomes a divisibility b·c + μ ∣ (a quadratic expression in the remaining unknowns), and the divisor moves with c. The step lemma shows that from a fibre point with positive modulus $t = b·c + μ$ one can always find a prime p, congruent to 1 modulo F·b·t and larger than it, and a new fibre point whose modulus is exactly t·p (\texttt{moving\_\allowbreak{}modulus\_\allowbreak{}prime\_\allowbreak{}step\_\allowbreak{}data}). The congruence is what allows the quadratic residue condition to be transported to the new modulus; the size condition is what makes the sequence strictly increasing. Iterating from any base point with positive modulus gives the moving-modulus step orbit, with the modulus factor F pinned to 2|det G₀| in the singular-block branch — G₀ being the nondegenerate quotient block of Lemma 2, not the original Gram matrix, whose determinant is zero in that branch — and to 8·max(1,|D'|)·|2k| in the rank-one coupled branch. Dirichlet's theorem on primes in arithmetic progressions (mathlib \texttt{Nat.\allowbreak{}infinite\_\allowbreak{}setOf\_\allowbreak{}prime\_\allowbreak{}and\_\allowbreak{}eq\_\allowbreak{}mod}) supplies the prime at each step, and only there.

\emph{Bounded branches.} Where no family is produced, the branch yields an explicit bound of the shape fboundG K ℓ q, and the constants are combined by taking the maximum over the finitely many branches. The branch table, with the kind emitted by each branch, ships with the formalization.

\end{proof}
